% Independent preview; original geometry and numeric relationships retained.
% 教学构造：横纵坐标均为归一化特征，右图与左图使用完全相同的分裂。
\begin{tikzpicture}[x=1cm,y=1cm,font=\small]
  \fill[FigureBlue!9] (0,0) rectangle (1.8,3.6);
  \fill[FigureRed!13] (1.8,0) rectangle (3.6,1.98);
  \fill[FigureBlue!9] (1.8,1.98) rectangle (3.6,3.6);
  \draw[book arrow] (0,0) -- (3.9,0) node[right] {$x_1$};
  \draw[book arrow] (0,0) -- (0,3.9) node[above] {$x_2$};
  \draw[FigureModel,thick] (1.8,0) -- (1.8,3.6);
  \draw[FigureModel,thick] (1.8,1.98) -- (3.6,1.98);
  \node[below] at (1.8,0) {$0.5$};
  \node[left] at (0,1.98) {$0.55$};
  \draw[FigureGrid,densely dashed] (0,1.98) -- (1.8,1.98);
  \foreach \x/\y in {0.4/0.5,0.8/1.2,1.3/0.8,0.6/2.4,1.4/3.1,2.3/2.5,3.1/3.1}
    \fill[FigureBlue] (\x,\y) circle (2pt);
  \foreach \x/\y in {2.2/0.6,2.8/1.3,3.2/0.5}
    \fill[FigureRed] (\x-0.065,\y-0.065) rectangle ++(0.13,0.13);
  \node[book node,minimum height=8mm,inner sep=4pt] (root) at (7.4,3.4) {$x_1\leq0.5$};
  \node[book node,fill=FigureBlue!9,minimum height=8mm,inner sep=4pt] (left) at (5.7,2.05) {类1};
  \node[book node,minimum height=8mm,inner sep=4pt] (right) at (8.6,2.05) {$x_2\leq0.55$};
  \node[book node,fill=FigureRed!13,minimum height=8mm,inner sep=4pt] (bottom) at (7.5,0.55) {类2};
  \node[book node,fill=FigureBlue!9,minimum height=8mm,inner sep=4pt] (top) at (9.6,0.55) {类1};
  \draw[book arrow] (root) -- node[above left] {是} (left);
  \draw[book arrow] (root) -- node[above right] {否} (right);
  \draw[book arrow] (right) -- node[above left] {是} (bottom);
  \draw[book arrow] (right) -- node[above right] {否} (top);
  \node[font=\figurefont\small] at (1.8,-0.7) {特征空间中的区域};
  \node[font=\figurefont\small] at (7.6,-0.7) {对应的二叉树};
\end{tikzpicture}
